\documentclass[11pt,a4paper]{article}

\usepackage[utf8]{inputenc}
\usepackage[T1]{fontenc}
\usepackage{amsmath,amssymb,amsthm}
\usepackage{geometry}
\usepackage{xcolor}
\usepackage{mdframed}
\usepackage{enumitem}
\usepackage{listings}
\usepackage{hyperref}
\geometry{margin=2.5cm}

\hypersetup{colorlinks=true, linkcolor=blue!60!black, urlcolor=blue!60!black}

\definecolor{boxbg}{RGB}{245,247,250}
\definecolor{boxline}{RGB}{90,120,170}
\definecolor{warnbg}{RGB}{253,245,230}
\definecolor{warnline}{RGB}{200,120,40}

\newmdenv[backgroundcolor=boxbg, linecolor=boxline, linewidth=1pt, roundcorner=4pt, skipabove=8pt, skipbelow=8pt, innertopmargin=8pt, innerbottommargin=8pt]{intuisi}
\newmdenv[backgroundcolor=warnbg, linecolor=warnline, linewidth=1pt, roundcorner=4pt, skipabove=8pt, skipbelow=8pt, innertopmargin=8pt, innerbottommargin=8pt]{catatan}

\lstset{basicstyle=\ttfamily\footnotesize, keywordstyle=\color{blue!60!black}\bfseries, commentstyle=\color{green!40!black}\itshape, stringstyle=\color{red!60!black}, numbers=left, numberstyle=\tiny\color{gray}, numbersep=8pt, frame=single, rulecolor=\color{gray!40}, breaklines=true, showstringspaces=false, language=Python, backgroundcolor=\color{gray!5}}

\newcommand{\E}{\mathbb{E}}
\newcommand{\clip}{\operatorname{clip}}

\title{\textbf{Modul Belajar: Safe Reinforcement Learning\\ dari Policy Gradient sampai PPO-Lagrangian}}
\author{Catatan studi --- MountainCar-v0}
\date{\today}

\begin{document}
\maketitle
\tableofcontents
\bigskip
\hrule
\bigskip

\section{Gambaran Besar}

Tujuan project: melatih sebuah \emph{agen} menyelesaikan game \textbf{MountainCar-v0},
namun \textbf{sambil mematuhi sebuah batasan keselamatan} (jangan ngebut melebihi
ambang kecepatan tertentu). Ini adalah masalah \emph{Safe Reinforcement Learning}:
memaksimalkan hadiah sekaligus menjaga sebuah biaya pelanggaran tetap di bawah batas.

\begin{intuisi}
\textbf{Dua tujuan yang saling bertarung:}
\begin{itemize}[nosep]
  \item \emph{Reward} --- cepat mencapai puncak bukit.
  \item \emph{Cost} --- jangan melampaui batas kecepatan.
\end{itemize}
Karena mobil \textbf{butuh} kecepatan untuk menanjak, kedua tujuan ini secara
intrinsik tegang. Di sinilah letak kesulitan (dan keindahan) masalahnya.
\end{intuisi}

\section{Fondasi: Markov Decision Process}

RL dimodelkan sebagai \textbf{Markov Decision Process} (MDP), tuple
$(\mathcal{S}, \mathcal{A}, P, r, \gamma)$:
\begin{itemize}[nosep]
  \item $\mathcal{S}$ --- himpunan \textbf{state} (mis. posisi \& kecepatan mobil).
  \item $\mathcal{A}$ --- himpunan \textbf{action} (mis. dorong kiri / diam / kanan).
  \item $P(s'\mid s,a)$ --- probabilitas transisi ke state $s'$ setelah aksi $a$.
  \item $r(s,a)$ --- \textbf{reward} yang diterima.
  \item $\gamma \in [0,1)$ --- \textbf{discount factor}.
\end{itemize}

Loop interaksinya berulang:
\[ \text{lihat } s_t \;\longrightarrow\; \text{pilih } a_t \;\longrightarrow\; \text{terima } r_t,\; s_{t+1} \;\longrightarrow\; \dots \]

\subsection{Return dan Discount}
Yang dimaksimalkan bukan reward sesaat, melainkan \textbf{return} (total reward terdiskon):
\[ G_t \;=\; \sum_{k=0}^{\infty} \gamma^{k}\, r_{t+k}. \]
Faktor $\gamma^k$ membuat reward jauh di masa depan bernilai lebih kecil, sekaligus
menjaga jumlahnya tetap berhingga.

\begin{intuisi}
Kadang aksi yang reward-nya jelek \emph{sekarang} membuka jalan ke reward besar
\emph{nanti} --- persis mobil yang harus mundur dulu sebelum bisa menanjak.
Itu sebabnya kita memaksimalkan \emph{return}, bukan reward sesaat.
\end{intuisi}

\section{Policy dan Policy Gradient}
\textbf{Policy} adalah strategi agen: sebuah distribusi peluang atas aksi,
\[ \pi_\theta(a \mid s), \]
diparametrikan oleh $\theta$ (bobot jaringan neural). Belajar $=$ menyetel $\theta$.

Objektif yang dimaksimalkan:
\[ J(\theta) \;=\; \E_{\tau \sim \pi_\theta}\!\left[\sum_{t} \gamma^t r_t\right], \]
dengan $\tau$ sebuah \emph{trajectory} (satu jalannya game).

\subsection{Teorema Policy Gradient}
\[ \boxed{\;\nabla_\theta J(\theta) \;=\; \E\!\left[\, \nabla_\theta \log \pi_\theta(a_t \mid s_t)\; A_t \,\right]\;} \]
\begin{itemize}[nosep]
  \item $\nabla_\theta \log \pi_\theta(a_t\mid s_t)$ --- arah yang \textbf{menaikkan peluang} aksi $a_t$.
  \item $A_t$ --- \textbf{advantage}, bobot yang menentukan seberapa kuat aksi itu didorong.
\end{itemize}

\begin{intuisi}
Aturannya satu kalimat: \emph{kalau suatu aksi ternyata lebih bagus dari perkiraan
($A_t>0$), perbesar peluangnya; kalau lebih jelek ($A_t<0$), perkecil.}
Agen tidak butuh "jawaban benar" --- ia membandingkan kenyataan dengan ekspektasinya sendiri.
\end{intuisi}

\section{Value Function dan Advantage}
\textbf{Value function} $V_\phi(s)$ (si \emph{critic}, jaringan neural kedua)
menebak return mulai dari state $s$:
\[ V(s) \;=\; \E\!\left[\, G_t \mid s_t = s \,\right]. \]
\textbf{Advantage} mengukur seberapa lebih bagus suatu aksi dibanding rata-rata di state itu:
\[ A_t \;=\; Q(s_t,a_t) - V(s_t). \]
Dalam praktik dipakai \textbf{TD error} (estimasi 1-langkah):
\[ \delta_t \;=\; r_t + \gamma V(s_{t+1}) - V(s_t), \]
yaitu (reward nyata $+$ nilai state berikutnya) $-$ (tebakan nilai sekarang) --- sebuah "kejutan".

\subsection{Generalized Advantage Estimation (GAE)}
PPO menghaluskan advantage sebagai rata-rata berbobot dari TD error:
\[ A_t^{\mathrm{GAE}} \;=\; \sum_{l=0}^{\infty} (\gamma \lambda_{\mathrm{gae}})^{l}\, \delta_{t+l}, \]
dengan $\lambda_{\mathrm{gae}}\in[0,1]$ menyeimbangkan \emph{bias} vs \emph{variance}.
(Catatan: $\lambda_{\mathrm{gae}}$ ini berbeda dari $\lambda$ Lagrangian di Bagian~\ref{sec:safe}.)

\section{PPO: Proximal Policy Optimization}
Policy gradient mentah rapuh: \textbf{satu update terlalu besar bisa merusak policy
secara permanen}. PPO mengatasinya dengan prinsip \emph{tetap dekat} dengan policy lama.

\subsection{Probability Ratio}
\[ r_t(\theta) \;=\; \frac{\pi_\theta(a_t\mid s_t)}{\pi_{\theta_{\mathrm{old}}}(a_t\mid s_t)} \]
\begin{itemize}[nosep]
  \item $r_t = 1$ --- policy tidak berubah.
  \item $r_t = 1.2$ --- policy baru 20\% lebih sering memilih aksi tersebut.
\end{itemize}

\subsection{Clipped Surrogate Objective}
\[ \boxed{\; L^{\mathrm{CLIP}}(\theta) \;=\; \E\!\left[\, \min\!\Big( r_t(\theta)\,A_t,\;\; \clip\big(r_t(\theta),\,1-\epsilon,\,1+\epsilon\big)\,A_t \Big) \right] \;} \]
dengan $\epsilon$ (mis. $0.2$) lebar batas. Mekanisme \texttt{min}$+$\texttt{clip}:
\begin{itemize}[nosep]
  \item \textbf{Jika $A_t > 0$} (aksi bagus): suku ter-clip memberi \emph{plafon} $1.2\,A_t$. Begitu $r_t$ melewati $1.2$, objektif berhenti naik $\Rightarrow$ gradien nol $\Rightarrow$ tak ada insentif melompat lebih jauh.
  \item \textbf{Jika $A_t < 0$} (aksi jelek): \texttt{min} memilih suku yang lebih pesimis, koreksi tetap diberikan namun terkendali.
\end{itemize}

\begin{intuisi}
Seperti memperbaiki resep dari feedback pelanggan dengan aturan
\emph{"maksimal ubah satu sendok bumbu per percobaan"}: pelan, tapi stabil dan tak
pernah merusak resep dalam semalam. Itulah seluruh rahasia kestabilan PPO.
\end{intuisi}

\subsection{Loss Value Function}
Critic dilatih dengan kuadrat selisih terhadap return nyata $R_t$:
\[ L^{V}(\phi) \;=\; \E\!\left[\big( V_\phi(s_t) - R_t \big)^2\right]. \]

\subsection{Langkah Update: Di Mana \texttt{lr} Muncul}
Perhatikan: \texttt{lr} (learning rate) \textbf{tidak muncul} di rumus objektif mana pun
($L^{\mathrm{CLIP}}$, $L^V$, teorema policy gradient). Itu wajar. Persamaan-persamaan
tersebut hanya mendefinisikan \emph{apa} yang dioptimalkan dan menghasilkan \emph{arah}
(gradien). \texttt{lr} baru masuk di langkah terpisah: \emph{bagaimana} mengambil langkahnya.
\[ \theta \;\leftarrow\; \theta + \alpha\,\nabla_\theta L^{\mathrm{CLIP}}(\theta), \qquad \phi \;\leftarrow\; \phi - \alpha_V\,\nabla_\phi L^{V}(\phi). \]
\begin{itemize}[nosep]
  \item $\nabla_\theta L$ --- menentukan \textbf{arah} (ke mana naik).
  \item $\alpha$ (=\texttt{lr}) --- menentukan \textbf{panjang langkah}.
\end{itemize}

\begin{intuisi}
Gradien itu \textbf{kompas} (menunjuk arah naik); \texttt{lr} itu \textbf{panjang langkahmu}.
Objektif cuma menggambar petanya; \texttt{lr} adalah keputusan seberapa jauh melangkah tiap
kali --- terpisah dari peta. \texttt{lr} besar = cepat tapi mudah overshoot; \texttt{lr}
kecil = stabil tapi lambat.
\end{intuisi}

Bandingkan dengan update $\lambda$ Lagrangian, yang aturannya ditulis \emph{eksplisit},
sehingga laju belajarnya ($\eta=$\texttt{lam\_lr}) \textbf{memang terlihat} di persamaan:
\[ \log\lambda \;\leftarrow\; \log\lambda + \eta\,(J_C - d). \]
Untuk $\theta$ dan $\phi$, langkah update biasanya disembunyikan di balik optimizer
(Adam/SGD) --- itu sebabnya \texttt{lr}-nya tak tampak di rumus objektif. Catatan: Adam
juga memodifikasi langkah efektif per-parameter, jadi \texttt{lr} yang diset bukan persis
panjang langkah akhir, tapi tetap berperan sebagai pengali besar-kecilnya update.

\subsection{Satu Epoch PPO}
\begin{enumerate}[nosep]
  \item Jalankan policy sekarang, kumpulkan $N$ langkah pengalaman (mis. \texttt{steps\_per\_epoch}=4000).
  \item Hitung advantage $A_t$ (via GAE).
  \item Naikkan $L^{\mathrm{CLIP}}$ terhadap $\theta$ (dengan rem clipping).
  \item Turunkan $L^{V}$ terhadap $\phi$.
  \item Ulangi sebanyak \texttt{epochs}.
\end{enumerate}

\section{Safe PPO: PPO-Lagrangian}
\label{sec:safe}
Sekarang ada \textbf{dua} sinyal: reward menghasilkan $A_t^{R}$, dan cost (pelanggaran
kecepatan) menghasilkan $A_t^{C}$. Masalahnya menjadi optimisasi berkendala:
\[ \max_\theta \; J_R(\theta) \quad \text{s.t.} \quad J_C(\theta) \le d, \]
dengan $J_C$ ekspektasi cost dan $d$ batasnya (\texttt{cost\_limit}$=15$).

\subsection{Lagrangian}
Bentuk Lagrangian-nya memperkenalkan pengali $\lambda \ge 0$:
\[ \mathcal{L}(\theta,\lambda) \;=\; J_R(\theta) - \lambda\big(J_C(\theta) - d\big). \]
Untuk policy, advantage gabungan yang dimasukkan ke rumus PPO menjadi:
\[ \boxed{\; A_t \;=\; A_t^{R} - \lambda\, A_t^{C} \;} \]

\begin{intuisi}
$\lambda$ adalah \textbf{"kenop tarif denda"} untuk ngebut. Jika $\lambda$ kecil,
suku $\lambda A_t^{C}$ nyaris nol $\Rightarrow$ agen mengabaikan cost $\Rightarrow$ ngebut terus.
Jika $\lambda$ besar, agen takut melanggar.
\end{intuisi}

\subsection{Update Dual untuk $\lambda$}
$\lambda$ dinaikkan lewat \textbf{gradient ascent} pada masalah dual:
\[ \boxed{\; \lambda \;\leftarrow\; \big[\,\lambda + \eta\,(J_C - d)\,\big]_+ \;} \]
dengan $\eta=$\texttt{lam\_lr} dan $[\cdot]_+$ proyeksi ke $\lambda\ge 0$.
Selama melanggar ($J_C > d$), $\lambda$ naik $\Rightarrow$ denda makin mahal,
sampai $J_C$ mendarat di $d$.

\section{Studi Kasus: Kenapa $\lambda$ Tidak Tumbuh}
\subsection{Gejala}
\begin{itemize}[nosep]
  \item avg cost: $26$--$37$ (target $\le 15$) --- constraint dilanggar.
  \item $\lambda$ final: $0.027$--$0.040$ --- sangat kecil, bahkan turun dari init $\approx 0.05$.
  \item Seed 1 return collapse ke $11$ (seed lain $\sim 44$--$48$).
\end{itemize}

\subsection{Diagnosa: Self-Suppression}
Implementasi menggunakan parametrisasi $\log\lambda$ dengan loss:
\begin{lstlisting}
loss_lambda = -torch.exp(log_lambda) * (actual_batch_cost - cost_limit)
\end{lstlisting}
Gradiennya terhadap $\log\lambda$:
\[ \frac{\partial\,\text{loss}}{\partial \log\lambda} = -\exp(\log\lambda)\,(J_C - d) = -\lambda\,(J_C - d). \]
Faktor $\lambda \approx 0.05$ \textbf{ikut mengalikan setiap update}. Langkah efektif di ruang log:
\[ \Delta \log\lambda \;\approx\; \eta\,\lambda\,(J_C - d) \;\approx\; 0.005 \times 0.05 \times 15 \;\approx\; 0.004 \;\text{per epoch.} \]

\begin{catatan}
\textbf{Self-suppression:} justru saat $\lambda$ kecil, ia menahan pertumbuhannya
sendiri. Di epoch awal (mobil belum ngebut, $J_C < d$) $\lambda$ malah didorong
\emph{turun}, lalu terjebak di nilai sangat kecil --- persis pola $\lambda$ final
yang lebih kecil dari nilai inisial.
\end{catatan}

\subsection{Penyebab Collapse Seed 1}
MountainCar butuh $|v|$ hingga $\sim 0.07$ untuk mencapai goal, sedangkan cost
menghukum $|v| > 0.04$. Constraint dan task \textbf{saling bertentangan}. Saat
$\lambda$ secara lokal menang, agen mengorbankan goal demi menurunkan kecepatan
$\Rightarrow$ return jatuh. Ini tension Pareto cost-vs-reward, diperparah osilasi Lagrangian.

\subsection{Perbaikan Minimal}
\paragraph{(a) Hilangkan faktor \texttt{exp}, normalisasi violation, tambah clamp:}
\begin{lstlisting}
with torch.no_grad():
    violation = (actual_batch_cost - cost_limit) / cost_limit  # ternormalisasi
    log_lambda += lam_lr * violation       # tanpa faktor exp -> no self-suppression
    log_lambda.clamp_(min=-2.0, max=3.0)   # floor & ceiling
lam = torch.exp(log_lambda)
\end{lstlisting}

\paragraph{(b) Hyperparameter:}
\[ \texttt{lam\_lr}:\; 0.005 \rightarrow 0.03, \qquad \texttt{log\_lambda\_init}:\; -3.0 \rightarrow 0.0 \;(\lambda\approx 1.0). \]

\paragraph{Alasan:} mulai dari $\lambda$ yang sebanding skala reward agar constraint
langsung "menggigit"; \texttt{lam\_lr} cukup besar agar responsif dalam 150 epoch
tanpa overshoot liar; normalisasi violation memisahkan \texttt{lam\_lr} dari skala cost mentah.

\begin{catatan}
\textbf{Peringatan jujur:} karena ambang $0.04$ melawan kebutuhan velocity, target
avg cost $\le 15$ mungkin menuntut pengorbanan sebagian return. Kalau setelah fix
return ikut turun di semua seed, itu sinyal \texttt{cost\_limit}/ambang yang perlu
dilonggarkan --- bukan lagi soal tuning $\lambda$.
\end{catatan}

\section{Ringkasan Hyperparameter}
Tiga kategori besaran --- jangan tertukar: \emph{hyperparameter} kamu set sendiri,
\emph{parameter} disetel training, \emph{besaran} dihitung dari data.

\begin{center}
\begin{tabular}{lll}
\hline
\textbf{Kategori} & \textbf{Contoh} & \textbf{Asalnya} \\
\hline
Hyperparameter (PPO)    & $\gamma$, \texttt{lr}, $\epsilon$, \texttt{lam\_gae}, \texttt{ent\_coef} & kamu ketik di config \\
Hyperparameter (safety) & \texttt{cost\_limit}, \texttt{lam\_lr}, \texttt{log\_lambda\_init} & kamu ketik di config \\
Parameter dipelajari    & $\theta$ (policy), $\phi$ (critic), $\lambda$ (Lagrangian) & disetel saat training \\
Besaran dihitung        & $\tau$, $\delta_t$, $A_t$, $r_t$, $G_t$ & dihitung dari data \\
\hline
\end{tabular}
\end{center}

\bigskip
Aturan cepat membedakan: kalau \textbf{dihitung dari data} ($A_t,\delta_t,r_t,\tau,G_t$)
$\Rightarrow$ bukan hyperparameter. Kalau \textbf{disetel jaringan saat training}
($\theta,\phi$) $\Rightarrow$ parameter yang dipelajari. Kalau \textbf{kamu ketik angkanya
di config sebelum run} $\Rightarrow$ baru itu hyperparameter.

\section{Ringkasan Simbol}
\begin{center}
\begin{tabular}{ll}
\hline
\textbf{Simbol} & \textbf{Makna} \\
\hline
$\pi_\theta(a\mid s)$ & policy: peluang aksi $a$ di state $s$ \\
$\theta,\ \phi$       & bobot jaringan policy (actor) \& value (critic) \\
$G_t,\ R_t$           & return (total reward terdiskon) \\
$\gamma$              & discount factor \\
$V(s),\ A_t$          & value function \& advantage \\
$\delta_t$            & TD error \\
$r_t(\theta)$         & probability ratio (policy baru / lama) \\
$\epsilon$            & lebar clipping PPO \\
$\alpha,\ \alpha_V$   & learning rate policy \& critic (\texttt{lr}) \\
$\lambda$             & pengali Lagrangian (denda cost) \\
$\lambda_{\mathrm{gae}}$ & parameter GAE (bias--variance) \\
$J_C,\ d$             & ekspektasi cost \& batasnya (\texttt{cost\_limit}) \\
$\eta$                & laju belajar $\lambda$ (\texttt{lam\_lr}) \\
\hline
\end{tabular}
\end{center}

\end{document}
